% 替换位置：摘要中“针对问题二”至“针对问题四”的三段；保留其余摘要
% 这是修改片段，不是完整文档；不要整份直接追加到原文末尾。
针对问题二，建立货箱组批、多点访问、异构机队与共享电池联合调度模型，
以医疗和首批保障时限为硬约束，采用多起点构造与模拟退火邻域搜索求解。
方案使用8架运输无人机和14组电池，通过24架次完成80箱按期交付，
加权延误为0，最晚返航时间为$103.478\,\mathrm{min}$，
总能耗为$69.241\,\mathrm{kWh}$。
本文候选池中，两种目标优先级选中相同方案。

针对问题三，基于三维轨迹、双向链路预算和地形遮挡判定，
将中继服务时间窗口纳入运输排序与机型选择。
方案由24个运输架次和3个中继架次完成，使用2架中继机和3组能源组件，
联合返航时间为$105.037\,\mathrm{min}$，总能耗为$72.435\,\mathrm{kWh}$。
独立回放覆盖爬升、巡航、下降和交接阶段，
经时空分辨率加密，在所述分辨率下未发现通信中断，
但尚未形成连续通信的数学证明。

针对问题四，冻结第三问基准调度，
在组内资源可重新编号、跨组中继完整复制的工作假设下，
将15个服务区合并为8个不可拆分块，利用区间峰值与区间着色核算最小资源，
完整枚举127个两组分区和966个三组分区。
两组方案缺少1架C型运输机和1架中继机；
三组方案缺少1架B型运输机、1架C型运输机、1组B型电池和2架中继机。
两组在资源缺口、配置规模及复制能耗方面更有利。
独立分区的资源需求由组内占用峰值及跨组中继任务复制共同决定。


% 替换位置：2.2中第二段（以“为此，在统一能耗模型的基础上”开头）
% 这是修改片段，不是完整文档；不要整份直接追加到原文末尾。
为此，在统一能耗模型的基础上，建立逐航段载荷更新、交付完成时间和资源占用模型，
将无人机与电池作为两类独立资源调度。
在多起点构造的可行解基础上，采用模拟退火邻域搜索，
通过机型调整、批次重排、移箱、换箱、合批、拆批及双批次重组所得方案，
以加权延误、全部运输返航时间、能耗和架次数为字典序目标。
最后通过独立确定性事件回放核验货箱完整性、硬时限、设备冲突和满电复用条件。


% 替换位置：第6.5节起至第6章结束；整体替换原6.5--6.8
% 这是修改片段，不是完整文档；不要整份直接追加到原文末尾。
\subsection{多起点构造与模拟退火邻域搜索}
\subsubsection{可行方案构造与资源解码}
首先利用多起点随机扰动构造完整可行方案，再将每个架次的货箱集合、机型及批次顺序作为搜索状态。
随机性仅用于搜索，不表示需求、天气或设备状态具有随机变化。
对每个候选架次枚举所涉及服务区的访问顺序，并重新检查质量、体积、逐段载荷、能耗和硬时限。
同型号机体与电池分别维护可用时间，架次最早准备时刻为
\begin{equation}
s_k^{\mathrm{earliest}}=\max\{r_{u_k}^{\mathrm{body}},r_{e_k}^{\mathrm{battery}}\}.
\end{equation}
机体返航后恢复可用，电池返航后充电，只有充满才允许下一次准备。
候选方案改变后重新解码后续资源与任务时刻，不沿用已经失效的资源安排。

\subsubsection{邻域与候选接受}
邻域操作包括机型变更、批次插入或交换、货箱移动或交换、批次合并和拆分，
以及两个批次合并后重新分配货箱或按目的地重组。
以完成时间和延误为主的辅助评价函数驱动模拟退火，允许暂时接受部分较差候选以离开局部解；
最终保留方案仍按式~\eqref{eq:q2_objective}的字典序目标筛选。
不满足硬约束的候选不进入接受比较，不能用惩罚项放宽医疗和首批截止时间。

\subsubsection{搜索设置与范围}
邻域搜索采用20260925、20260926和20260927三个随机种子，每组搜索时间预算为$180\,\mathrm{s}$。
每架次访问点数上限取3或4，通过不同搜索设置构造候选方案。
按照字典序目标筛选满足第二问约束的候选，并对选定方案独立核验。
访问点数上限、邻域和壁钟时间均为计算限制，不是原题约束；未穷尽全部可行调度。

\subsection{调度结果与分析}
\subsubsection{基准方案汇总}
所得方案完成80箱按期交付。16箱医疗物资与30箱首批保障物资均满足各自硬时限，
两类标记允许重叠。结果见表~\ref{tab:q2_summary}。
\begin{table}[htbp]
\centering
\small
\caption{第二问基准调度结果}
\label{tab:q2_summary}
\begin{tabular}{lr}
\toprule
指标 & 结果 \\
\midrule
交付箱数 & 80 \\
加权延误$D$ & 0 \\
运输架次数 & 24 \\
多服务区架次数 & 6 \\
最晚返航时刻/s & 6208.684865 \\
总能耗/kWh & 69.240573590 \\
实际运输机/电池数量 & 8/14 \\
最低返航SOC & $20.958111\%$ \\
最小硬时限余量/s & 194.866445 \\
\bottomrule
\end{tabular}
\end{table}
A、B、C型分别执行12、6、6架次。各实体机均执行3架次，U08的末次返航决定$C_{\max}$。
该方案的最晚返航时间为$103.478081\,\mathrm{min}$，全部运输任务在给定库存内完成。
第一问的资源与时限约束、路线限制和目标顺序均与本问不同，不能仅比较18与24架次判断优劣。
\begin{table}[htbp]
\centering
\small
\caption{所得方案的实体无人机任务分配}
\label{tab:q2_units}
\begin{tabular}{llrrr}
\toprule
无人机 & 机型 & 架次数 & 累计占用/s & 末次返航/s \\
\midrule
U01 & A & 3 & 5390.093 & 5797.832 \\
U02 & A & 3 & 5699.660 & 5817.252 \\
U03 & A & 3 & 5215.296 & 5215.296 \\
U04 & A & 3 & 5337.313 & 6113.782 \\
U05 & B & 3 & 6024.529 & 6024.529 \\
U06 & B & 3 & 5418.168 & 5418.168 \\
U07 & C & 3 & 6195.854 & 6195.854 \\
U08 & C & 3 & 6208.685 & 6208.685 \\
\bottomrule
\end{tabular}
\end{table}

\subsubsection{多点架次的载荷与时间示例}
选取C型无人机U08执行的P009架次，其使用电池C-BAT-03，携带8箱物资，路线为
\begin{equation}
\mathrm{O01}\rightarrow\mathrm{S006}\rightarrow\mathrm{S007}\rightarrow\mathrm{O01}.
\end{equation}
初始载荷$74\,\mathrm{kg}$，在S006交付4箱后降至$37\,\mathrm{kg}$，
在S007交付其余4箱后空载返回。
\begin{table}[htbp]
\centering
\small
\caption{P009架次的逐航段载荷与能耗}
\label{tab:q2_multistop}
\begin{tabular}{lrrrr}
\toprule
航段 & 载荷/kg & 起始/s & 到达/s & 能耗/kWh \\
\midrule
O01$\rightarrow$S006 & 74 & 1968.159 & 2301.527 & 2.011156 \\
S006$\rightarrow$S007 & 37 & 2625.527 & 3023.843 & 1.040271 \\
S007$\rightarrow$O01 & 0 & 3347.843 & 3907.628 & 1.438274 \\
\bottomrule
\end{tabular}
\end{table}

该架次准备开始与起飞时刻分别为$1428.158932\,\mathrm{s}$和$1968.158932\,\mathrm{s}$，
两站交付完成时刻为$2625.527260\,\mathrm{s}$和$3347.843322\,\mathrm{s}$，
均在到达之后计入完整交接时间。
全架次能耗$4.489701013\,\mathrm{kWh}$，返航SOC为$43.878737\%$。
机体于$3907.628219\,\mathrm{s}$返航，但电池到$5956.922243\,\mathrm{s}$才充满，
表明机体与能源恢复可用时间必须分别处理。

\subsection{目标优先级对比与最优性边界}
\subsubsection{目标优先级比较}
在本文搜索候选池内，分别按$D\rightarrow C_{\max}\rightarrow E\rightarrow N$
和$C_{\max}\rightarrow D\rightarrow E\rightarrow N$筛选。
以$10^{-7}\,\mathrm{s}$作为完成时间的数值比较容差，两者选中相同方案。
\begin{table}[htbp]
\centering
\small
\caption{第二问候选池的目标优先级比较}
\label{tab:q2_priority}
\begin{tabular}{lrrrr}
\toprule
排序 & 加权延误 & 完成时间/s & 能耗/kWh & 架次 \\
\midrule
$D\to C_{\max}\to E\to N$ & 0 & 6208.684865 & 69.240573590 & 24 \\
$C_{\max}\to D\to E\to N$ & 0 & 6208.684865 & 69.240573590 & 24 \\
\bottomrule
\end{tabular}
\end{table}

两种排序得到相同结果仅反映所搜索范围，不证明及时性与完成时间在全部可行方案中没有冲突。

\subsubsection{最优性边界}
由于$D\geq0$且当前方案达到$D=0$，第一层目标达到其下界；
后续完成时间、能耗与架次数没有全局最优证明。
因此本问应表述为独立验证通过的当前搜索最好可行方案，不能仅凭多轮随机搜索声称全局最优。

\subsection{独立事件回放与可行性检查}
独立程序重新读取原始货箱、设备和时限，采用40位十进制精度复算能耗与时刻，
并按准备开始、返航充电、充满恢复可用三类事件进行回放，共72个运输资源事件。
检查80箱无遗漏、无重复、无错投，质量、体积、逐段载荷、SOC及全部硬时限均合规，
机体占用和电池使用、充电、复用无冲突。
最低返航SOC为$20.958111\%$，最小硬时限余量$194.866445\,\mathrm{s}$，
最大架次能耗重算差约$4.595\times10^{-16}\,\mathrm{kWh}$。
时间、能耗、质量、体积与SOC容差分别为$10^{-7}\,\mathrm{s}$、$10^{-8}\,\mathrm{kWh}$、
$10^{-9}\,\mathrm{kg}$、$10^{-12}\,\mathrm{m^3}$及$10^{-10}$。
判定使用未舍入数据，同刻先释放后占用。
全部24架次、80箱交付和电池充电记录另附；第二问的运输验证不能替代第三问的通信检查。


% 替换位置：第7.4节起至第7章结束；整体替换原7.4--7.6
% 这是修改片段，不是完整文档；不要整份直接追加到原文末尾。
\subsection{候选生成与联合搜索方法}
\subsubsection{空间候选与服务结构}
采用水平网格、服务区位置和地形高点生成中继候选，离地高度取150米和300米，
经DEM范围、高度及回传链路筛选后形成597个候选位置。
时间候选组合涉及其中49个位置和1904组三架次结构。
求解采用候选筛选与通信服务日历相结合的方法，检查中继服务窗口与运输轨迹的时间匹配。

\subsubsection{通信服务日历参与运输调度}
进一步采用一架中继持续服务、另一架先后执行两个中继任务的结构，
将服务窗口、返航、机体周转、再次准备与建链时间纳入运输候选的可执行性检查。
以运输可行候选作为搜索起点，调整运输机型、批次顺序、访问顺序及准备时刻，
使需要中继的航段避开中继尚未建链或换场的时段。
三处位置沿用候选130、511和489，最终按实际通信需求重新确定服务起止时刻。
运输时刻与中继服务窗口共同参与候选筛选。

\subsubsection{计算限制}
服务日历搜索种子为20260928，时间预算180秒，包含最终细检查的实际耗时约182.748秒。
候选运输批次最多涉及4个服务区；中继位置与三架次结构仍受限制。
本研究服务日历搜索暂要求所有已给定期望时限均满足，以寻找$D=0$方案，
这只是搜索限制，普通物资期望时间在原模型中仍为软目标。
未穷举连续位置、任意组批和中继任务数量，也未证明全部离散候选范围内的最优性。

\subsection{联合调度结果与权衡分析}
\subsubsection{中继位置与任务安排}
所得方案使用2架中继机、3组初始满电能源组件执行3个中继架次。
悬停位置、任务时刻和能耗分别列于下表，所有位置均满足DEM范围及离地高度约束。
\begin{table}[htbp]
\centering
\small
\caption{所得方案的中继悬停位置}
\label{tab:q3_locations}
\begin{tabular}{llrrr}
\toprule
架次 & 机体 & 经度/度 & 纬度/度 & 海拔/m \\
\midrule
R001 & R01 & 109.204167 & 23.048056 & 837.268860 \\
R002 & R02 & 109.270134 & 23.015271 & 627.188568 \\
R003 & R02 & 109.266111 & 23.056389 & 623.466492 \\
\bottomrule
\end{tabular}
\end{table}
\begin{table}[htbp]
\centering
\small
\caption{所得方案的中继任务安排}
\label{tab:q3_relay_schedule}
\begin{tabular}{lrrrrr}
\toprule
架次 & 准备/s & 建链完成/s & 服务结束/s & 返航/s & 能耗/kWh \\
\midrule
R001 & 325.265 & 1056.880 & 5721.485 & 6302.231 & 1.717471 \\
R002 & 130.908 & 738.785 & 3329.431 & 3768.932 & 1.019583 \\
R003 & 4315.907 & 5077.631 & 5506.123 & 6099.161 & 0.457479 \\
\bottomrule
\end{tabular}
\end{table}

R001、R002、R003的返航SOC分别为$46.329029\%$、$68.138039\%$和$85.703779\%$。
R002于$3768.932416\,\mathrm{s}$返航，完成300秒周转后于$4068.932416\,\mathrm{s}$恢复机体可用，
早于R003的准备时刻$4315.906989\,\mathrm{s}$，两次机体占用不冲突。
但R002的能源组件到$4683.137907\,\mathrm{s}$才充满，不能直接承接R003。

\subsubsection{联合方案汇总}
80箱全部按期交付，医疗及首批硬时限全部满足，$D=0$。
\begin{table}[htbp]
\centering
\small
\caption{第三问联合调度结果}
\label{tab:q3_summary}
\begin{tabular}{lr}
\toprule
指标 & 结果 \\
\midrule
运输/中继架次 & 24/3 \\
运输最晚返航/s & 6208.684865 \\
联合最晚返航/s & 6302.230771 \\
运输能耗/kWh & 69.240573590 \\
中继能耗/kWh & 3.194532893 \\
总能耗/kWh & 72.435106483 \\
运输机/运输电池 & 8/14 \\
中继机/能源组件 & 2/3 \\
\bottomrule
\end{tabular}
\end{table}

联合完成时间为$105.037180\,\mathrm{min}$，由R001的最终返航决定，
比运输最晚返航晚$93.545906\,\mathrm{s}$。
第二问与第三问选定的运输子方案具有相同的路线、机型和任务时刻。
该运输子方案满足第二问的运输约束；第三问还需配置相应中继任务，
并独立检查全过程通信，才能形成联合可行方案。

\subsubsection{目标优先级比较}
按$D\rightarrow C_{\mathrm{joint}}\rightarrow E_{\mathrm{total}}\rightarrow N_{\mathrm{total}}$
与$C_{\mathrm{joint}}\rightarrow D\rightarrow E_{\mathrm{total}}\rightarrow N_{\mathrm{total}}$
筛选当前候选池，均选中上述$D=0$、$6302.230771\,\mathrm{s}$、
$72.435106483\,\mathrm{kWh}$及24运输加3中继架次的方案。
这不证明两个完整优化问题具有相同的全局最优解，尤其本研究服务日历搜索偏重全部按期交付的候选。

\subsection{独立轨迹回放与通信验证}
\subsubsection{任务与资源检查}
独立程序重读原始货箱、设备、通信参数和DEM，重构运输全过程轨迹，复算中继飞行与能耗，
检查所有机体、运输电池及中继组件的占用、周转和充电。
两种优先级选中的方案均通过检查，运输最低返航SOC为$20.958111\%$，
最小硬时限余量为$194.866445\,\mathrm{s}$。

\subsubsection{分辨率加密检查}
检查覆盖爬升、巡航、下降、交接以及中继建链与服务边界，结果见表~\ref{tab:q3_refinement}。
\begin{table}[htbp]
\centering
\small
\caption{所得方案的通信分辨率核验}
\label{tab:q3_refinement}
\begin{tabular}{lrrr}
\toprule
时间步长与空间检查 & 点数 & 失败点 & 最小裕量/dB \\
\midrule
1秒，5米射线 & 36218 & 0 & 0.000703269 \\
0.5秒，2.5米射线 & 72113 & 0 & 0.000356918 \\
0.5秒及局部加密，全穿越像元 & 73220 & 0 & 0.000155292 \\
\bottomrule
\end{tabular}
\end{table}

对端点裕量低于$0.02\,\mathrm{dB}$或保障源变化的区间细化至不超过0.05秒，
对定位到的切换点二分至$10^{-7}\,\mathrm{s}$。
最差所选链路裕量出现在P010巡航阶段$t=3860.254957\,\mathrm{s}$附近。

\subsubsection{验证结论与适用范围}
准确结论为：\textbf{在所述时空分辨率下未发现中断}。
离散检查未严格排除采样间隙中的短暂失效，不能写成已数学证明中断时长为零。
全部196个运输阶段保守列入尚未严格认证区间；248条通信保障记录及其任务关联用于第四问冻结。
第一层$D=0$达到非负下界，后续目标无全局最优证明。
标准航程能量标定、爬升效率解释、未计运输交接悬停能耗，
以及中继多机并发与零额外切换耗时仍是需要用户核验的工作假设。


% 替换位置：8.3中“实际计算得到14个不可拆分服务区块”所在段；其余变量与公式保留
% 这是修改片段，不是完整文档；不要整份直接追加到原文末尾。
实际计算得到8个不可拆分服务区块：
\begin{align*}
\mathcal B_1&=\{\mathrm{S001,S002,S009,S010,S012}\},&
\mathcal B_2&=\{\mathrm{S003,S006,S007}\},\\
\mathcal B_3&=\{\mathrm{S004}\},&
\mathcal B_4&=\{\mathrm{S005,S011}\},\\
\mathcal B_5&=\{\mathrm{S008}\},&
\mathcal B_6&=\{\mathrm{S013}\},\\
\mathcal B_7&=\{\mathrm{S014}\},&
\mathcal B_8&=\{\mathrm{S015}\}.
\end{align*}
这些块由第三问运输架次的共同访问关系求传递闭包得到，分区决策作用于上述不可拆分块。


% 替换位置：8.4.2从“未分区情况下重新计算得到最小配置”起到8.5之前
% 这是修改片段，不是完整文档；不要整份直接追加到原文末尾。
未分区且资源可共享时，重新计算得到最小配置
\begin{equation}
\boldsymbol R^0=(4,2,2,6,4,4,2,3),
\end{equation}
向量顺序为A、B、C型运输机，A、B、C型电池，中继机和中继能源组件。
在$4315.906989\,\mathrm{s}$，R003开始预留能源，R001仍在使用其组件，
R002组件尚在充电，三组能源占用同时存在，因此至少需要3组中继组件。
R002组件到$4683.137907\,\mathrm{s}$才充满，晚于R003准备开始，不能复用承接。
因此，该冻结时序下的中继能源组件最小需求为3组，与具体可用时刻的核算结果一致。


% 替换位置：第8.6节起至第8章结束；整体替换原8.6--8.11；8.5目标公式保留
% 这是修改片段，不是完整文档；不要整份直接追加到原文末尾。
\subsection{求解方法与最优性范围}
对8个不可拆分块进行无标签非空分区枚举，分区数分别为
\begin{equation}
\left\{\begin{matrix}8\\2\end{matrix}\right\}=127,
\qquad
\left\{\begin{matrix}8\\3\end{matrix}\right\}=966.
\end{equation}
缓存$2^8-1=255$个非空子集的资源需求和工作量后计算各分区目标。
以$12G$和$12S$进行整数比较，均衡层比较$\mathrm{CV}^2$，并列容差$10^{-12}$。
全部分区均被检查，两个基准分区在该容差下均唯一。
最优性仅针对冻结第三问基准、完整中继复制、组内可重新编号及所采用目标，
不能扩展为运输、通信与分区共同决策的全局最优性。

\subsection{分区与配置结果}
\subsubsection{两组分区}
记$\mathcal V$为全部15个服务区。两组最优分区为
\begin{align}
\mathcal V_1&=\mathcal V\setminus\{\mathrm{S008}\},&
\mathcal V_2&=\{\mathrm{S008}\},\\
\boldsymbol n_1&=(4,2,2,6,4,3,2,3),&
\boldsymbol n_2&=(0,0,1,0,0,1,1,1).
\end{align}
总需求为$\boldsymbol R^{(2)}=(4,2,3,6,4,4,3,4)$，
缺少1架C型运输机和1架中继机；中继能源组件剩余2组。
目标为$G=1$、$S=8.666667$、$\mathrm{CV}=0.735372$。

\subsubsection{三组分区}
三组最优分区为
\begin{align}
\mathcal V_1&=\mathcal V\setminus\{\mathrm{S008,S015}\},\\
\mathcal V_2&=\{\mathrm{S008}\},\qquad \mathcal V_3=\{\mathrm{S015}\}.
\end{align}
配置分别为
\begin{align}
\boldsymbol n_1&=(4,2,2,6,4,3,2,3),\\
\boldsymbol n_2&=(0,0,1,0,0,1,1,1),\\
\boldsymbol n_3&=(0,1,0,0,1,0,1,1).
\end{align}
总需求$\boldsymbol R^{(3)}=(4,3,3,6,5,4,4,5)$，
缺少1架B型运输机、1架C型运输机、1组B型电池和2架中继机；组件剩余1组。
目标为$G=2.25$、$S=10.083333$、$\mathrm{CV}=0.914171$。
\begin{table}[htbp]
\centering
\small
\caption{两种基准分区的资源需求与缺口}
\label{tab:q4_resources}
\begin{tabular}{lrrrrr}
\toprule
类型 & 库存 & 两组需求 & 两组缺口 & 三组需求 & 三组缺口 \\
\midrule
A型运输机 & 4 & 4 & 0 & 4 & 0 \\
B型运输机 & 2 & 2 & 0 & 3 & 1 \\
C型运输机 & 2 & 3 & 1 & 3 & 1 \\
A型电池 & 6 & 6 & 0 & 6 & 0 \\
B型电池 & 4 & 4 & 0 & 5 & 1 \\
C型电池 & 4 & 4 & 0 & 4 & 0 \\
中继机 & 2 & 3 & 1 & 4 & 2 \\
中继组件 & 6 & 4 & 0 & 5 & 0 \\
\bottomrule
\end{tabular}
\end{table}
\begin{table}[htbp]
\centering
\small
\caption{分区的各组工作量}
\label{tab:q4_workload}
\begin{tabular}{llrrrrrr}
\toprule
$K$ & 组 & 箱数 & 质量/kg & 体积/m$^3$ & 运输架次 & 中继任务 & 工作量/s \\
\midrule
2 & G1 & 75 & 713 & 1.882 & 23 & 3 & 55588.142 \\
2 & G2 & 5 & 45 & 0.129 & 1 & 1 & 8476.668 \\
3 & G1 & 72 & 688 & 1.815 & 22 & 3 & 53759.898 \\
3 & G2 & 5 & 45 & 0.129 & 1 & 1 & 8476.668 \\
3 & G3 & 3 & 25 & 0.067 & 1 & 1 & 8105.209 \\
\bottomrule
\end{tabular}
\end{table}

\subsection{资源缺口与额外配置的来源}
\subsubsection{组间错峰复用损失}
两组中，大组在$t=0$由P003与P005同时占用两架C型机，S008组的P012另需一架，
组间禁止调配导致总需求3架。大组由R001、R002重叠占用两架中继机，
S008组另需R001完整副本，因此中继需求也是3架。
三组再独立配置S015组，需要额外保有B型机、B型电池及中继机。
各组峰值不必同刻出现；独立执行要求求各组峰值之和，而不是分组后的全局峰值。

\subsubsection{库存剩余、额外配置与任务重复}
库存剩余为$U_\gamma=\max(0,I_\gamma-R_\gamma)$，
额外配置为$A_\gamma=R_\gamma-R^0_\gamma$，
完整中继复制引起的新增任务与能耗另行统计，三者不能混称冗余。
两组相较未分区增加1架C型机、1架中继机及1组组件；
三组增加1架B型机、1架C型机、1组B型电池、2架中继机和2组组件。
增加组件配置不等于组件库存短缺。

两组对R001额外复制一次，三组对R001额外复制两次，故中继任务总数分别为4次与5次。
复制能耗分别为
\begin{align}
E_{\mathrm{copy}}^{(2)}&=1.717471065\,\mathrm{kWh},\\
E_{\mathrm{copy}}^{(3)}&=3.434942131\,\mathrm{kWh}.
\end{align}
对应新增机体占用分别为$6276.966104\,\mathrm{s}$和$12553.932208\,\mathrm{s}$，
新增能源组件占用分别为$7174.688724\,\mathrm{s}$和$14349.377449\,\mathrm{s}$。
完整复制保持物理任务时刻，故两种分区仍有$D=0$、
$C_{\mathrm{joint}}=6302.230771\,\mathrm{s}$；运输与中继总能耗分别为
\begin{align}
E_{\mathrm{total}}^{(2)}&=74.152577548\,\mathrm{kWh},\\
E_{\mathrm{total}}^{(3)}&=75.870048613\,\mathrm{kWh}.
\end{align}

\subsection{目标优先级与解释口径比较}
\subsubsection{配置规模与工作量均衡}
将$G\to S\to\mathrm{CV}$改为$G\to\mathrm{CV}\to S$，两种$K$均选中相同基准分区。
若将均衡置于首位，则可降低工作量差异，但资源缺口明显增加。
\begin{table}[htbp]
\centering
\small
\caption{冻结方案的目标优先级比较}
\label{tab:q4_priority}
\begin{tabular}{clrrr}
\toprule
$K$ & 排序 & $G$ & $S$ & $\mathrm{CV}$ \\
\midrule
2 & $G\to S\to\mathrm{CV}$ & 1.000000 & 8.666667 & 0.735372 \\
2 & $G\to\mathrm{CV}\to S$ & 1.000000 & 8.666667 & 0.735372 \\
2 & $\mathrm{CV}\to G\to S$ & 4.500000 & 12.500000 & 0.005234 \\
3 & $G\to S\to\mathrm{CV}$ & 2.250000 & 10.083333 & 0.914171 \\
3 & $G\to\mathrm{CV}\to S$ & 2.250000 & 10.083333 & 0.914171 \\
3 & $\mathrm{CV}\to G\to S$ & 7.333333 & 15.333333 & 0.107693 \\
\bottomrule
\end{tabular}
\end{table}

\subsubsection{禁止复制中继架次}
若禁止复制，则同一中继实际保障的运输任务也必须同组。
将这些绑定加入运输关系并求传递闭包后，15个服务区合并为一个不可拆分集合。
例如R001和R002共同关联包含S001、S002、S009、S010、S012的运输块，
R003关联的S004、S005也已与R001关联；三类中继任务的绑定连通，覆盖全部服务区。
因此该解释下不存在$K=2$或$K=3$的合法非空分区，
这是禁止复制解释下的分区不可行，不是原题一般意义上的无解。

\subsubsection{继承原实体编号}
按每组出现的不同原编号计数，并对跨组同名编号分别配置实体，不能跨组共享。
本文固定基准分区的继承编号计数，与组内重新编号的最小资源向量一致；
在继承编号口径下重新优化分区，也得到相同基准需求。
这仅是本文冻结时序的结果，不说明重新编号在一般情况下没有作用。

\subsection{独立事件回放与结果有效性}
独立程序重读原始参数，复算货箱交付、能耗、SOC、充电和周转，并回放组内资源事件。
15个服务区无遗漏或重复，同架次访问点同组，80箱及其交付时刻保持不变，
16箱医疗物资和30箱首批保障物资硬时限全部满足，最小余量$194.866445\,\mathrm{s}$。
248条实际通信保障区间被继承，中继副本完整，组内资源无冲突且不跨组共享。
共12个导出场景通过检查，另以独立枚举核对127个两组分区与966个三组分区，得到10项目标核对证书。
冻结输入SHA-256在检查前后保持一致。
时间、能耗、SOC、质量、体积容差分别为$10^{-7}\,\mathrm{s}$、$10^{-8}\,\mathrm{kWh}$、
$10^{-10}$、$10^{-9}\,\mathrm{kg}$、$10^{-12}\,\mathrm{m^3}$。
通过检查表示所列配置可以承接冻结任务，不表示现有库存足够；两种$K$均无零缺口分区。
通信只继承第三问的数值验证范围，相同中继副本未额外认证共址干扰与实际网络容量。

\subsection{结果分析}
冻结方案下，两组的资源缺口、配置规模、工作量变异系数和复制能耗均低于三组，
因此按既定字典序目标优先采用两组。
分区资源代价主要来自独立保有设备后失去跨组错峰复用机会，
以及同一中继任务被多个组使用时所需的完整任务副本。
上述结论以第三问的任务时序和实际通信保障关系为固定输入，
依赖完整复制、组内同型号资源可重新编号等工作假设。
在此口径下，现有库存均不能独立承接两种分区，需要按资源类型补足缺口。


% 替换位置：第9章“这里是模型评价”占位句；以下为与本研究结果有关的评价段落
% 这是修改片段，不是完整文档；不要整份直接追加到原文末尾。
本模型统一四问的地形、载荷、能耗和资源占用规则，并用独立回放检查计算结果，
有助于避免逐段载荷、满电复用和跨问继承口径不一致。
批次重组、机型选择和通信服务日历共同参与求解，
第二问通过24个运输架次完成80箱按期交付，
第三问配置3个中继架次完成全过程数值通信保障。

模型也存在限制。第二、三问尚未穷尽货箱分配、运输顺序和连续中继位置，
零延误仅证明第一层目标达到下界，不能证明完整字典序全局最优。
通信检查仍为有限分辨率的数值验证，不能严格排除采样间隙失效。
能量标定、爬升效率、未计运输交接悬停耗能及无线并发容量等工作假设需要进一步核验。

此外，分区要求资源不跨组调配，因而会产生额外配置及中继任务复制能耗。
第四问的最优性仅针对冻结时序、完整中继复制和组内重新编号口径，
不能扩展为运输、通信与分区共同决策时的全局最优性。


% 替换位置：删除现有无关MATLAB最短路及Python演示MIP附录，用真实代码说明替代；放在appendices环境中
% 这是修改片段，不是完整文档；不要整份直接追加到原文末尾。
\section{求解代码与独立验证}
随结果附件提供实际运行的Python程序。
第二问的\texttt{search.py}完成批次与机型邻域搜索；
第三问的\texttt{joint.py}和\texttt{calendar\_joint.py}完成候选匹配与服务日历搜索；
第四问的\texttt{solve\_q4.py}进行无标签分区枚举和区间着色。
独立检查由\texttt{verify\_q2.py}、\texttt{verify\_joint.py}及\texttt{verify\_q4.py}执行，
\texttt{replay\_all.py}提供统一复算入口。
随机种子、搜索时限、依赖与复算命令见代码包说明，所有结论对应保存的未舍入结果。
